\par\bigskip
\begin{center}
\rule{.88\textwidth}{.3pt}\par
Göttingische gelehrte Anzeigen.\quad 24 February 1823.\par
\rule{.88\textwidth}{.3pt}
\end{center}
\pdfbookmark[1]{24 February 1823}{anzeige-1823-02-24}

A lecture presented to the Royal Society on 2 February by Councillor \textsc{Gauss}, entitled
\begin{center}
\textit{Theoria combinationis observationum erroribus minimis obnoxiae, pars posterior,}
\end{center}
\noindent is immediately connected with an earlier one, reviewed in these pages [26 February 1821]. We need only briefly recall that its purpose was to establish a new foundation for the so-called method of least squares. On this foundation the method appears as the most suitable way of combining observations, not approximately but with mathematical exactness; not only in the restricted case of a very large number of observations, and not depending on a hypothetical law for the probability of observational errors, but in complete generality. The present second part of the investigation now develops this doctrine further in a series of propositions and problems closely connected with it. It would not suit the arrange\SourceLeaf{106}%
ment of these pages to follow these investigations here step by step; nor is that necessary, since the treatise itself is already in press. We shall therefore confine ourselves to mentioning the subjects of some investigations that can more readily be isolated.

The values of the unknown quantities that accord with the method of least squares, and which may be called the \textit{safest values}, are found by means of a determinate elimination; the weights to be assigned to these determinations are found by an indeterminate elimination, as is already known from the \textit{Theoria motus Corporum Coelestium}. Here it is proved \textit{a priori}, in a new way, that under the assumptions in force this elimination is always possible. A remarkable symmetry among the coefficients arising in the indeterminate elimination is also demonstrated.

Although these elimination operations are, in general, easy and clear to survey, it cannot be denied that their actual numerical execution becomes burdensome when there is a considerable number of unknown quantities. As for the \textit{determinate} elimination, which suffices to find the safest values of the unknown quantities, the author has already indicated a procedure that shortens the actual calculation as much as the nature of the subject permits, in the \textit{Theoria Motus Corporum Coelestium}, and has developed it in detail in a paper in the first volume of the \textit{Commentt. Rec. Soc. R. Gott.}, entitled \textit{Disquisitio de elementis ellipticis Palladis}. This procedure also has the advantage that the weight of the determination of the one unknown quantity treated as the last in the calculation follows automatically. Since the order of the unknown quantities is entirely arbitrary, and any one of them may therefore be treated as last, this procedure suffices in every case where the weight is required for only one of the unknown quantities; the burdensome indeterminate elimination is then avoided.

The practice, which has since become so widespread among computational astronomers, of applying the method of least squares to difficult astronomical calculations—for example, to the complete determination of cometary orbits, where the number of unknown quantities rises to six—has, however, made keenly felt the need for finding the weights of the safest values of \textit{all} the unknown quantities \SourceLeaf{107}%
in a more convenient way than by indeterminate elimination.
Since the efforts of some geometers\footnote{For example, those of \textsc{Plana}. See \textit{Zeitschrift für Astronomie und verwandte Wissenschaften}, vol. 6, p. 258.} had not succeeded, the only recourse had been to carry out the algorithm mentioned above as many times as there were unknown quantities, changing their order each time and assigning each of them the last place once. It seems to us, however, that, compared with indeterminate elimination, this artless procedure gains nothing in brevity of calculation. The author has therefore made this important subject the object of a separate investigation and communicated a new algorithm for determining the weights of the values of \textit{all} the unknown quantities, which seems to possess all the flexibility and brevity that the nature of the subject can permit.

The safest value of a quantity that is a given function of the unknown quantities in the problem is found by substituting for the latter their safest values obtained by the method of least squares. A problem not hitherto treated, however, is how to find the weight to be assigned to that determination. The solution given here deserves all the more attention from computational astronomers, since it turns out that several of them have previously proceeded incorrectly in this matter.

As is well known, the sum of the squares of the differences between the quantities directly observed and those values obtained for them, as functions of the unknown quantities, by substituting the safest values for the latter, forms an absolute minimum. (If the observations have unequal reliability, each square must first be multiplied by its respective weight before the sum is taken.) Hence, as soon as one of the unknown quantities is assigned a value different from its safest value, a similar aggregate will always be larger than that minimum, whatever values may be assigned to the other unknown quantities. But the other unknown quantities can be determined in only one way so that the increase in the aggregate is as small as possible, or so that the aggregate itself becomes a relative minimum. The investigation carried out here by the au\SourceLeaf{108}%
thor leads to some interesting truths that shed still broader light on the whole doctrine.
It sometimes happens that only after an extensive calculation on a series of observations has been completed in every part does one learn of a new observation that one would like to include as well. In many cases it may be desirable not to have to begin the entire elimination process again from the start, but to be able to determine the modification caused by adding the new observation to the safest values and their weights. The author has therefore treated this problem separately, as well as the related problem in which one finds reason, after the fact, to assign a different weight from the one originally given to an observation already used, and wishes to obtain the changes in the final results without repeating the calculation from the beginning.

In a separate paper in the \textit{Zeitschrift für Astronomie und verwandte Wissenschaften} [March and April 1816], the author showed how the \textit{probable} error of a class of observations—the measure of their uncertainty used until now—can be found approximately from a sufficient number of actual observational errors. This procedure, like the use of probable error in general, depends, and had to depend, on the hypothetical form of the probability assigned to individual errors. The first part of the present treatise does show how, from the same data, the mean error of the observations—a suitable measure of their inaccuracy—can be found approximately. Yet the difficulty remains that, with full precision, one can seldom or almost never know the true magnitude of a set of actual observational errors. In practice, people have therefore always used the differences between what the observations yielded and the results of calculation based on the safest values of the unknown quantities on which the observations depend, as found by the method of least squares. But since one is not entitled to regard the safest values as the true values themselves, it is easy to see that this procedure must always make both the probable error and the mean error \textit{too small}, and therefore assign greater accuracy to the observations and the results de\SourceLeaf{109}%
rived from them than they really possess. Of course, in the case where the number of observations is many times greater than the number of unknown quantities, this inaccuracy is of little consequence. Yet the dignity of science requires, on the one hand, that we see fully and precisely how much error we risk making in this way; on the other, results of calculation have in fact often been presented in important cases by that defective procedure, where the stated assumption did not hold. The author has therefore made this subject the object of a separate investigation, which has led to a very remarkable and extremely simple result. To convert the mean error found by the indicated defective procedure into the correct one, it need only be multiplied by
\[
\surd\frac{\pi-\rho}{\pi}
\]
where \(\pi\) denotes the number of observations and \(\rho\) the number of unknown quantities.

The final investigation concerns the determination of the degree of accuracy to be assigned to this determination of the mean error itself; its results must, however, be read in the treatise itself.
\par\bigskip
\begin{center}
\rule{.88\textwidth}{.3pt}\par
\end{center}
